\section{Introduction}
\label{sec:intro}

Retrieval-augmented generation (RAG) answers a question from passages that a retrieval method returns.
Retrieval metrics such as Hit@$k$ and nDCG@10 score an intermediate outcome: the ranking of passages.
The final output of a RAG system is an answer, which a reader model produces from the retrieved passages.
A gain in a retrieval metric therefore need not translate proportionally to answer accuracy.
Retrieval methods also differ in computational cost, which is a separate dimension.

Existing studies measure different parts of this relationship.
MTEB evaluates text embedding models on 8 tasks and 58 datasets, one of which is retrieval \citep{muennighoff2022mteb}.
It evaluates retrieval performance directly rather than the effect of retrieval on a fixed downstream reader.
\citet{baidya2026choosing} compares a commercial embedding model with open-source models on four BEIR datasets and reports query latency and per-query cost.
That report states that it ``does not measure end-to-end RAG answer quality''.
RAGAS \citep{es2023ragas} defines reference-free metrics for faithfulness, answer relevance and context relevance.
Its evaluation does not compare retrievers.
ARES \citep{saadfalcon2024ares} scores RAG systems built from three retrievers and three generators with fine-tuned judges, and each system retrieves one passage.
ARES does not isolate the contribution of retrieval to answer quality and does not measure computational cost.
We do not claim that a gap between retrieval gains and answer gains is new.
Our focus is its controlled measurement, together with cost.

We ask four research questions.
\begin{itemize}
\item \textbf{RQ1:} How do open embedding models differ in retrieval quality on one fixed corpus?
\item \textbf{RQ2:} How much of a retrieval gain transfers to end-to-end answer accuracy (EM)?
\item \textbf{RQ3:} What does each model cost (encoding time, latency, index size, memory) on free hardware?
\item \textbf{RQ4:} Which models are Pareto-optimal on quality vs cost?
\end{itemize}

In this paper, we make four contributions.
\begin{itemize}
\item \textbf{(C1) Controlled isolation.} We vary only the retrieval method: BM25 and six open-source dense encoders. The reader (Qwen2.5-7B-Instruct-AWQ), the corpus (a 300,000-passage BEIR-NQ sub-corpus) and the 2,255 joined queries stay fixed (\S\ref{sec:method}).
\item \textbf{(C2) Transfer decomposition.} We report the ratio $\Delta\mathrm{EM}/\Delta\mathrm{Hit@5}$ against all-MiniLM-L6-v2 as the reference, with paired percentile bootstrap confidence intervals over queries and a pre-registered reporting rule (\S\ref{sec:method}).
\item \textbf{(C3) Measured cost on free hardware.} We measure corpus encoding time, query latency, index size and peak GPU memory on one free Tesla T4 (\S\ref{sec:results}).
\item \textbf{(C4) Pareto analysis.} We compute the quality--cost Pareto frontier from the measured costs. Quality comparisons along it use paired exact McNemar or Wilcoxon signed-rank tests with Holm correction over 21 pairs. A post-hoc sensitivity analysis treats cost differences of at most 5\% as ties (\S\ref{sec:results}).
\end{itemize}
We are not aware of work that measures retrieval-to-answer transfer ratios jointly with measured cost on free hardware.

Within our setup, $\sim$20--25\% of Hit@5 gain appears as EM gain (\S\ref{sec:results}).
Section~\ref{sec:related} reviews related work, Sections~\ref{sec:method}--\ref{sec:setup} describe the method and setup, Section~\ref{sec:results} reports results, and Sections~\ref{sec:discussion}--\ref{sec:conclusion} discuss them and state limits.
